\documentclass[12pt,oneside,a4paper,sumario=tradicional,portuguese]{abntex2}

% - PACOTES UTILIZADOS - %
\usepackage{lmodern} % Usa a fonte Latin Modern
\usepackage[T1]{fontenc} % Seleciona a codificacao de fonte T1, necessaria para o uso adequado de caracteres acentuados
\usepackage[utf8]{inputenc} % Permite a entrada de caracteres acentuados diretamente
\usepackage{lastpage} % Referencia a ultima pagina do documento
\usepackage{indentfirst} % Indenta o primeiro paragrafo de cada secao
\usepackage{color,colortbl} % Permite o uso de cores em texto e tabelas
\usepackage{graphicx} % Permite a inclusao de graficos e imagens
\usepackage{microtype} % Melhora a justificacao do texto
\usepackage[brazilian,hyperpageref]{backref} % Referencias bibliograficas com hiperlinks
\usepackage[alf]{abntex2cite} % Formata as citacoes e referencias bibliograficas conforme a ABNT
\usepackage{subcaption} % Permite subfiguras e sublegendas
\usepackage{amsmath} % Permite a formatacao avancada de formulas matematicas
\usepackage{amsfonts} % Permite o uso de fontes da American Mathematical Society
\usepackage{amssymb} % Permite o uso de simbolos matematicos adicionais
\usepackage{booktabs} % Melhora a formatacao de tabelas
\usepackage[table,xcdraw]{xcolor} % Permite cores em tabelas e desenhar tabelas coloridas
\usepackage{multirow} % Permite mesclar varias linhas em tabelas
\usepackage{rotating} % Permite a rotacao de tabelas e figuras
\usepackage{times} % Usa a fonte Times New Roman
\usepackage{float} % Controla a posicao dos elementos flutuantes
\usepackage[normalem]{ulem} % Permite sublinhado e outros formatos de texto
\useunder{\uline}{\ul}{} % Define o comando \uline para sublinhar texto
\usepackage{longtable} % Permite a criacao de tabelas que ocupam varias paginas
\usepackage[portuguese,ruled,lined]{algorithm2e} % Escreve algoritmos em portugues com linhas e regras
\usepackage{pdfpages} % Inclui paginas de documentos PDF
\usepackage{svg} % Permite a inclusao de graficos SVG
\svgsetup{inkscapelatex=false} % Configuracao para inclusao de SVGs sem usar LaTeX no Inkscape
\usepackage{tabularray} % Ferramentas avancadas para tabelas
\definecolor{Silver}{rgb}{0.753,0.753,0.753} % Define a cor Silver
\newcommand{\source}[1]{\caption*{Fonte: {#1}} } % Comando para adicionar fonte nas legendas
\newcommand{\SubItem}[1]{{\setlength\itemindent{15pt} \item[-] #1}} % Comando para criar subitens em listas
\newcommand{\mycomment}[1]{} % Comando para adicionar comentarios no codigo

% - INFORMACAO DO DOCUMENTO - %

\titulo{Ensino introdutório de \textit{machine learning} no Ensino Médio por meio de uma sequência didática apoiada por uma plataforma interativa}
\autor{Luiz Gabriel Bezerra da Conceição}
\orientador{Prof. Dr. Marcus de Barros Braga}
\local{Belém, PA}
\data{2026}
\instituicao{%
	Universidade Federal do Pará
	\par
	Instituto de Ciências Exatas e Naturais
	\par
	Programa de Pós-Graduação em Ensino de Computação para a Educação Básica (PROFCOMP)}
\tipotrabalho{Projeto de Qualificação de Mestrado}

\preambulo{Projeto de qualificação apresentado ao Programa de Pós-Graduação em Ensino de Computação para a Educação Básica da Universidade Federal do Pará, como requisito parcial para a qualificação no curso de Mestrado Profissional em Ensino de Computação.}

% - INFORMAÇÕES SOBRE O PDF - %

\makeatletter
\hypersetup{
		pdftitle={\@title}, 
		pdfauthor={\@author},
	    pdfsubject={\imprimirpreambulo},
	    pdfcreator={Luiz Gabriel Bezerra da Conceição},
		pdfkeywords={},
		colorlinks=true,% false: boxed links; true: colored links
	    linkcolor=black, % color of internal links
	    citecolor=blue, % color of links to bibliography
	    filecolor=magenta, % color of file links
		urlcolor=green,
		bookmarksdepth=4
}
\makeatother
\setlength{\parindent}{1.3cm}
\setlength{\parskip}{0.2cm}
\makeindex

% - COMANDOS DE SUBSTITUIÇÃO - %

\renewcommand{\backrefpagesname}{Citado na(s) página(s):~}
\renewcommand{\backref}{}
\renewcommand*{\backrefalt}[4]{
	\ifcase #1 %
		Nenhuma citação no texto.%
	\or
		Citado na página #2.%
	\else
		Citado #1 vezes nas páginas #2.%
	\fi}%
\renewcommand{\imprimirorientadorRotulo}{Orientador:}

\renewcommand{\imprimircapa}{
\begin{capa}
\center
\includesvg[width=0.2\textwidth]{imagens/logoUFPA.svg}

 { \ABNTEXchapterfont\large 
 UNIVERSIDADE FEDERAL DO PARÁ\\
 INSTITUTO DE CIÊNCIAS EXATAS E NATURAIS\\
 PROGRAMA DE PÓS-GRADUAÇÃO EM ENSINO DE COMPUTAÇÃO PARA A EDUCAÇÃO BÁSICA (PROFCOMP)\\}
   \vspace*{2.0cm}

 {\ABNTEXchapterfont\LARGE
 \textbf{\imprimirtitulo}}
 \vspace*{2cm}
 
  { \ABNTEXchapterfont\large
   \imprimirautor}   
   \vspace*{2.0cm}
 
 \vspace*{\fill}
 {\large\imprimirlocal}
 \par
 {\large\imprimirdata}
 \vspace*{1cm}
\end{capa}
}

\makeatletter
\renewcommand{\folhaderostocontent}{
\begin{center}
 { \ABNTEXchapterfont\large 
 UNIVERSIDADE FEDERAL DO PARÁ\\
 INSTITUTO DE CIÊNCIAS EXATAS E NATURAIS\\
 PROGRAMA DE PÓS-GRADUAÇÃO EM ENSINO DE COMPUTAÇÃO PARA A EDUCAÇÃO BÁSICA (PROFCOMP)\\}
   \vspace*{\fill}
   
    {\ABNTEXchapterfont\large~\imprimirautor}

    \vspace*{\fill}\vspace*{\fill}
    {\ABNTEXchapterfont\large
     \textbf{\imprimirtitulo}}
     \vspace*{\fill}

    \abntex@ifnotempty{\imprimirpreambulo}{
    \hspace{.45\textwidth}
    \begin{minipage}{.5\textwidth}
    	\imprimirpreambulo\\
    \end{minipage}%
    \vspace*{\fill}}
    
    {\ABNTEXchapterfont\large
    \imprimirorientadorRotulo
    \ABNTEXsectionfont~\imprimirorientador\par}

    \vspace*{\fill}
    {\large\imprimirlocal}
    \par
    {\large\imprimirdata}
    
    \vspace*{1cm}
\end{center}
}

% - DOCUMENTO PRINCIPAL - %

\begin{document}

% Retira espaço extra obsoleto entre as frases.
\frenchspacing 

% ELEMENTOS PRÉ-TEXTUAIS
\pretextual

% Capa
\imprimircapa

% Folha de Rosto
\imprimirfolhaderosto*

% Resumo
\setlength{\absparsep}{18pt} 

\begin{resumo}[Resumo]
 \begin{otherlanguage*}{brazil}
   
   \vspace{\onelineskip}
   \noindent 
   
Estudantes da Educação Básica convivem com aplicações de inteligência artificial, mas esse contato cotidiano está longe de assegurar a compreensão de como dados e regras orientam seu funcionamento. Este projeto de qualificação objetiva desenvolver e analisar uma sequência didática apoiada por uma plataforma interativa para o ensino introdutório de \textit{machine learning} no Ensino Médio. O tema entra como conteúdo de Computação, examinando como modelos computacionais partem de dados para gerar previsões numéricas, separar classes ou formar grupos. A proposta articula a Resolução CNE/CEB nº 1/2022 e o pensamento computacional ao uso de múltiplas representações, organizando as aulas da previsão inicial ao experimento prático para sustentar a observação e a explicação. O produto educacional une oito encontros a uma plataforma web executada localmente com dados sintéticos tabulares, cobrindo regressão linear, regressão logística, \textit{k}-vizinhos mais próximos (\textit{k}-NN), árvore de decisão e \textit{k}-means. De natureza aplicada e abordagem qualitativa apoiada por dados descritivos, a pesquisa ocorrerá como intervenção pedagógica em uma turma do próprio pesquisador. O corpus reunirá o diagnóstico inicial e o instrumento final com seu desafio de transferência, somados a três registros intermediários e às notas docentes. Uma rubrica examinará o que muda ou permanece nas dificuldades de relacionar os dados e a tarefa ao efeito de ajustes nas representações, considerando ainda os limites das decisões automatizadas. Frequências e transições entre níveis descreverão o conjunto, sem atribuir à plataforma um efeito causal isolado.
   
   \textbf{Palavras-chave}: Ensino de Computação; Machine Learning; Ensino Médio; Sequência Didática; Plataforma Interativa; Produto Educacional.

 \end{otherlanguage*}
\end{resumo}

% O resumo em inglês (Abstract) será incluído após a validação da versão portuguesa.


% Lista de figuras, tabelas e siglas
\hypersetup{pdfborder = {0 0 0}}

% Lista de Tabelas
\pdfbookmark[0]{\listtablename}{lot}
\renewcommand{\listtablename}{Lista de tabelas}
\listoftables*
\cleardoublepage
  
\begin{siglas} 
  \item [BNCC] 		 Base Nacional Comum Curricular
  \item [CEB] 		 Câmara de Educação Básica
  \item [CNE] 		 Conselho Nacional de Educação
  \item [CNS] 		 Conselho Nacional de Saúde
  \item [IA] 		 Inteligência Artificial
  \item [k-NN] 		 \textit{k}-Nearest Neighbors (\textit{k}-Vizinhos Mais Próximos)
  \item [MEC] 		 Ministério da Educação
  \item [ML] 		 \textit{Machine Learning} (Aprendizado de Máquina)
  \item [PPA] 		 Projeto Permanente por Afinidade
  \item [PROFCOMP] 	 Programa de Pós-Graduação em Ensino de Computação para a Educação Básica
  \item [SBC] 		 Sociedade Brasileira de Computação
  \item [UFPA] 	 	 Universidade Federal do Pará      
\end{siglas}


% Sumário
% inserir o sumario

\hypersetup{%
    pdfborder = {0 0 0}
}
\pdfbookmark[0]{\contentsname}{toc}
\tableofcontents*
\cleardoublepage

% Elementos textuais com numeração arábica
\pagenumbering{arabic}
% Reinicia a contagem do número de páginas
\setcounter{page}{1}


% ELEMENTOS TEXTUAIS
\textual

% - CAPÍTULOS - %
\chapter{Introdução}

\section{Contextualização e delimitação do tema}

A incorporação de recursos de inteligência artificial às rotinas escolares e sociais dos estudantes aproxima os jovens de respostas prontas que podem ser usadas para concluir tarefas sem que se compreenda o conteúdo envolvido. Esse movimento se intensifica nas redes sociais e em plataformas de vídeo, onde circulam demonstrações de ferramentas que realizam atividades acadêmicas e prometem resultados rápidos.

Assistir a essas demonstrações, porém, não garante a aprendizagem do conteúdo nem esclarece o funcionamento da inteligência artificial, deixando encobertos tanto os dados utilizados quanto o processo que gerou o resultado. Sem esse entendimento, o estudante pode ter dificuldade para analisar criticamente a resposta recebida, e o currículo escolar pode oferecer pouco espaço para explicar como um sistema de inteligência artificial chega a seus resultados. A Resolução CNE/CEB nº 1/2022 define diretrizes para a Computação na Educação Básica como complemento à Base Nacional Comum Curricular e atribui aos sistemas de ensino a organização da implantação curricular \cite{brasil2022computacao}. Em 2026, o Ministério da Educação publicou o documento orientador \textit{Inteligência artificial na Educação Básica}, que reforça o papel do professor e defende que a IA seja tratada como recurso de tutoria e apoio à aprendizagem, e não apenas como geradora de respostas prontas \cite{brasil2026iaeducacaobasica}.

A educação em inteligência artificial tem sido debatida internacionalmente como parte da formação cidadã, pois os estudantes precisam compreender, conforme sua escolaridade, as possibilidades e os limites de sistemas inteligentes \cite{touretzky2019envisioning}. No Ensino Médio, essa necessidade se coloca de modo particular: usar muito os recursos digitais não é o mesmo que ter letramento computacional. Muitos estudantes usam celulares e aplicativos com desenvoltura, mas tiveram poucas oportunidades de ler uma tabela com autonomia para distinguir correlação de previsão, percebendo na prática por que a alteração da regra muda o resultado. Outros ainda enfrentam restrições de acesso a computador e internet fora da escola.

A discussão sobre o uso responsável deve ser acompanhada pelo ensino dos conceitos que explicam como a inteligência artificial gera suas respostas. Se a conversa se limita a liberar ou proibir uma ferramenta, o estudante permanece na posição de usuário. Compreender como uma inteligência artificial é treinada pode ajudá-lo a entender como ela produz respostas e a utilizá-la de modo favorável ao seu desenvolvimento cognitivo. O ensino de Computação permite deslocar a conversa para o que alimenta um modelo e por que uma mudança nos dados altera sua resposta, e este projeto focaliza o ensino introdutório de \textit{machine learning} para discuti-lo a partir da relação entre dados e modelos computacionais.

O recorte dispensa conhecimentos prévios de programação ou a formação de especialistas em algoritmos. Em oito encontros, estudantes do Ensino Médio partirão da leitura de dados para testar hipóteses ao modificar parâmetros, usando a Matemática na interpretação dos resultados sem torná-la um obstáculo inicial.

\section{Problema educacional e justificativa}

Na rotina escolar, alguns estudantes recorrem à inteligência artificial para obter resultados imediatos, sem necessariamente se apropriar do conhecimento envolvido. Quando esse uso se resume a buscar respostas prontas, o funcionamento da ferramenta fica oculto para quem a utiliza. Nesse cenário, ensinar conceitos básicos de \textit{machine learning} ajuda o jovem a conhecer o processo por trás das respostas que recebe. Soma-se a isso a falta de oportunidades de trabalho com dados. Em geral, tabelas e gráficos apresentados no Ensino Médio chegam prontos, inseridos em questões. Nessas situações, a leitura dos valores pode ocorrer sem uma discussão sobre a origem dos dados ou sobre o que mudaria no resultado se um valor ou uma condição fosse diferente.

Compreender um modelo de \textit{machine learning} exige algum contato prévio com dados, além de noções básicas de estatística e probabilidade. Já existem trabalhos voltados ao ensino do tema na Educação Básica. Em um mapeamento de 30 unidades didáticas, \citeonline{marques2020teaching} observaram que várias delas cobrem apenas os processos mais acessíveis, como a gestão de dados, ou apresentam o aprendizado e o teste dos modelos de forma abstrata, deixando parte do funcionamento interno oculto. Isso sugere que, nessas unidades, o contato com dados pode se restringir à gestão, sem mostrar como os dados e os parâmetros alteram o resultado do modelo. No Ensino Médio, a revisão de \citeonline{martins2023findings} indica que os estudantes conseguem compreender conceitos básicos de \textit{machine learning} e aplicá-los em tarefas com algoritmos. A mesma revisão aponta a necessidade de avaliações mais rigorosas sobre esse ensino. A presença da inteligência artificial na escola também pede atividades que vinculem os dados e a finalidade da tarefa ao modo como regras e parâmetros alteram o resultado do modelo.

Quando os dados chegam prontos e as relações entre eles não são trabalhadas, o estudante tende a aprender os conceitos de \textit{machine learning} apenas pelo nome. O vocabulário básico da área, que vai do treinamento do modelo ao ajuste de parâmetros para gerar previsões, acaba repetido sem conexão real. Um estudante pode dizer que um modelo foi treinado sem saber explicar o que muda quando um parâmetro é alterado, e pode reconhecer uma reta no gráfico sem relacioná-la a uma previsão. O problema educacional está na passagem entre nomear um conceito e utilizá-lo para explicar uma situação.

Há ainda uma diferença importante entre apresentar a inteligência artificial como tema de cultura digital e ensiná-la como conteúdo de Computação. A discussão cultural é indispensável para problematizar impactos sociais, abrangendo desde vieses e riscos até questões de autoria. No ensino de Computação, o estudante também precisa investigar, na prática, uma representação de dados e uma regra computacional. A proposta não reduz questões éticas a uma equação ou a um gráfico. Considera-se que a discussão sobre limites passa a ter outra base quando o estudante reconhece que modelos resultam de decisões humanas de modelagem.

Romper com a recepção passiva de resultados prontos exige que o estudante faça previsões antes de alterar dados ou parâmetros para observar e explicar o efeito produzido. Esse processo começa antes do modelo, na origem dos dados. Um estudante que monta uma tabela com medidas da própria turma sabe de onde vêm os números e o que cada coluna representa. Essa experiência torna a leitura do gráfico menos mecânica. Em seguida, ver uma reta se deslocar quando os dados mudam, ou a mudança de categoria quando o limiar é alterado, produz evidências visíveis. Ao explicar o que observou, o estudante pode relacionar o efeito visual à alteração que realizou. Essa organização se aproxima das abordagens práticas, baseadas em problemas ou projetos, que \citeonline{martins2023findings} identificaram como bem-sucedidas no Ensino Médio. Compreender de onde vêm os dados e o que muda com um parâmetro é o que permite ao estudante questionar o resultado que recebe, e não apenas aceitá-lo.

O produto educacional prevê uma plataforma interativa em português brasileiro, articulada a uma sequência didática para estudantes do Ensino Médio. A plataforma apresenta cinco módulos introdutórios: regressão linear, regressão logística, \textit{k}-vizinhos mais próximos, árvore de decisão e \textit{k}-means. Os modelos serão tratados como conteúdos centrais, sem receber necessariamente a mesma formalização matemática. Com a regressão linear, a turma analisa a linha de tendência para fazer previsões e medir o erro, passando depois à probabilidade e ao limiar na regressão logística. Já o \textit{k}-NN destaca distância e votação, a árvore de decisão explicita regras e o \textit{k}-means agrupa dados não rotulados.

O uso da plataforma, por si só, não será tomado como evidência de aprendizagem. Ferramentas visuais podem apoiar a compreensão, mas seu potencial depende de como as representações são desenhadas e das tarefas cognitivas pedidas aos estudantes \cite{ainsworth2006deft}. Por isso, cada módulo partirá de uma situação-problema e da origem dos dados para que a turma levante uma hipótese antes da manipulação controlada. A atividade avança, assim, da previsão inicial ao experimento prático, culminando na explicação do que foi observado. A devolutiva do experimento participa da mediação docente, aproximando-se da avaliação formativa discutida por \citeonline{black1998assessment}. Serão oito encontros de cerca de 50 minutos. No segundo, a turma poderá organizar uma tabela simples antes de estudar regressão linear; nos demais módulos, trabalhará com cenários sintéticos fornecidos pelo site. Um estudante que modifica o limiar de uma classificação precisa discutir por que a mudança ocorreu. Sem essa discussão, a manipulação do controle pode ficar desvinculada do conceito que se pretende ensinar.

A justificativa profissional está na necessidade de um recurso aplicável por professores da Educação Básica em turmas heterogêneas, considerando a limitação de tempo e a variação de infraestrutura entre escolas. Em um programa profissional como o ProfComp, a relevância do estudo reside em articular um problema escolar ao conteúdo de Computação por meio da intervenção docente com um produto utilizável. A análise da plataforma considera seu uso em aula, pois o estudo descarta a ideia de que uma interface interativa isolada resolva dificuldades de aprendizagem.

Espera-se contribuir para aproximar a inteligência artificial do Ensino Médio regular, mantendo visível como cada modelo funciona sem exigir o aprofundamento próprio dos cursos técnicos. A contribuição é localizada e dependerá das condições de aplicação. Documentar essas condições pode ajudar outros docentes a avaliar se uma proposta semelhante serve ao seu contexto, sem antecipar generalizações que a pesquisa não poderá sustentar.

\section{Pergunta de pesquisa}

Que mudanças e permanências podem ser identificadas nas explicações de estudantes do Ensino Médio sobre fundamentos introdutórios de \textit{machine learning} ao longo de uma sequência didática apoiada por uma plataforma interativa?

A pergunta dirige a análise para o modo como os estudantes justificam escolhas e explicam gráficos ou tabelas em diferentes momentos da sequência. Em vez de medir o efeito isolado da plataforma, o estudo examinará o que muda ou permanece nas dificuldades de relacionar os dados e a tarefa ao ajuste de parâmetros na regra do modelo. As produções serão elaboradas ao longo de um percurso que combina o uso individual da plataforma à mediação docente, respeitando as condições reais da turma.


\chapter{Objetivos}

\section{Objetivo geral}

Desenvolver e analisar uma sequência didática apoiada por uma plataforma interativa para o ensino introdutório de fundamentos de \textit{machine learning}, como conteúdo de Computação, a estudantes do Ensino Médio.

\section{Objetivos específicos}

Os objetivos específicos desta pesquisa compreendem:

\begin{itemize}
    \item analisar referências sobre o ensino de inteligência artificial e \textit{machine learning} na Educação Básica, articulando o uso de representações visuais interativas à avaliação conceitual;
    \item relacionar os objetivos de aprendizagem da sequência aos eixos e às habilidades pertinentes da Computação na Educação Básica, sem supor que a norma prescreva algoritmos específicos;
    \item organizar uma sequência didática sobre regressão linear, regressão logística, \textit{k}-vizinhos mais próximos, árvore de decisão e \textit{k}-means, partindo de previsões e experimentos para observar e explicar resultados;
    \item ajustar às condições do campo de aplicação tanto os controles e explicações matemáticas da plataforma quanto os cenários sintéticos e tabelas organizadas pela turma;
    \item aplicar a sequência em uma turma de Ensino Médio em que o pesquisador é o professor;
    \item identificar o que muda ou permanece nas dificuldades dos estudantes ao explicarem como os dados e a tarefa se ligam a regras e parâmetros no resultado, confrontando o diagnóstico inicial com os registros de aula e o instrumento final;
    \item examinar como os estudantes mobilizam os conceitos estudados diante de uma situação-problema inédita em relação às aulas;
    \item sistematizar a plataforma junto à sequência didática e seus instrumentos como produto educacional.
\end{itemize}

\chapter{Referencial Teórico}

\section{Ensino de Computação, inteligência artificial e formação na Educação Básica}

O ensino de Computação na Educação Básica não se limita ao uso de ferramentas digitais. Ele envolve conhecimentos e práticas que permitem compreender o funcionamento de sistemas computacionais. Recursos cotidianos baseados em inteligência artificial aproximam essa discussão da experiência dos estudantes, mas essa proximidade não substitui o ensino dos princípios que orientam esses recursos.

A Resolução CNE/CEB nº 1/2022 estabelece normas para a Computação como complemento à BNCC e atribui aos sistemas de ensino a definição de parâmetros e abordagens para sua implementação \cite{brasil2022computacao}.

O ensino de Computação difere da informática na educação, em que o computador apoia a aprendizagem de conteúdos de diversas áreas e o foco recai sobre sua inserção pedagógica \cite{valente2002uso}. A plataforma deste projeto também será um recurso didático, mas o objeto central de aprendizagem não é operar a interface nem estudar um conteúdo externo com auxílio do computador.

A proposta focaliza conceitos e práticas da Computação relacionados ao funcionamento e aos limites de sistemas de inteligência artificial. Espera-se que os estudantes interpretem a participação de dados e regras na saída gerada, comparando a finalidade dos modelos ao antecipar e explicar cada modificação.

A Matemática apoia essa leitura com gráficos e funções, além de noções de distância, média, erro e probabilidade, sem substituir o foco computacional das atividades. A Computação na Educação Básica se estrutura pelos eixos Pensamento Computacional, Mundo Digital e Cultura Digital, conforme as Diretrizes da Sociedade Brasileira de Computação e a Resolução CNE/CEB nº 1/2022 \cite{sbc2018diretrizes,brasil2022computacao}. Integrar esses eixos evita reduzir a plataforma ao manuseio de controles técnicos. No Ensino Médio, a seleção de conteúdos segue as habilidades da norma que articulam análise de dados e modelos computacionais à inteligência artificial \cite{brasil2022computacao}.

O pensamento computacional participa da proposta na organização do problema a partir de seus dados e regras, longe de equivaler a qualquer uso do computador. \citeonline{wing2006computational} chama a atenção para práticas como abstração e decomposição na formulação de soluções computacionais. Neste estudo, essa aproximação aparece quando o estudante separa uma situação em dados disponíveis e regras do modelo. Essa mobilização apoia a leitura do conteúdo específico da sequência, afastando a pretensão de medir um desenvolvimento geral do pensamento computacional ou de atestar ganhos diretos em outras disciplinas.

No Mundo Digital, os estudantes analisam como modelos processam representações numéricas e executam regras para produzir respostas que podem orientar decisões automatizadas. A Cultura Digital aparece no exame do uso de sistemas de inteligência artificial presentes no cotidiano.

\citeonline{touretzky2019envisioning} defendem que a educação em inteligência artificial para a Educação Básica deve ajudar os estudantes a compreender ideias que orientam sistemas inteligentes e seus impactos. O documento da AI4K12 sobre \textit{machine learning} \cite{touretzky2023machine} situa a aprendizagem a partir de dados como uma das grandes ideias da inteligência artificial e inclui, para a faixa do Ensino Médio, a diferenciação entre aprendizagem supervisionada e não supervisionada. Essas referências não são um currículo pronto para a realidade brasileira, mas oferecem linguagem conceitual para selecionar conhecimentos que não dependam de programação avançada.

\citeonline{camada2020ensino} identificaram possibilidades de pesquisa e aplicação para o ensino de inteligência artificial na Educação Básica a partir de uma revisão sistemática. \citeonline{duraes2021ensino} relatam uma sequência didática sobre fundamentos de IA com estudantes do Ensino Médio integrado. Esses estudos mostram que o tema já vem sendo trabalhado na Educação Básica. Na elaboração desta proposta, a mediação docente precisa considerar tanto o repertório matemático da turma quanto o tempo e a infraestrutura disponíveis.

Neste projeto, a inteligência artificial permite discutir como sistemas computacionais dependem de decisões humanas de modelagem. A questão educacional não será se o estudante consegue usar uma ferramenta generativa para terminar uma tarefa escolar. O que se buscará verificar é se, em situações delimitadas, o estudante consegue explicar como alterações em parâmetros e regras afetam previsões e classificações, de modo que a proposta permaneça vinculada ao ensino de Computação.

\section{\textit{Machine learning} introdutório no Ensino Médio}

O \textit{machine learning} ocupa lugar específico entre os conteúdos de inteligência artificial porque permite examinar como computadores identificam padrões em dados e os utilizam para produzir respostas em tarefas definidas. Ao mapearem unidades instrucionais destinadas à escola, \citeonline{marques2020teaching} observaram que vários materiais apresentam apenas aspectos básicos de \textit{machine learning} ou deixam partes do funcionamento do modelo em nível abstrato.

A proposta dispensa o detalhamento técnico de cada algoritmo. O estudante precisa compreender o suficiente para reconhecer como dados e regras participam da produção dos resultados do modelo.

\citeonline{martins2023findings} reúnem evidências de que estudantes iniciantes do Ensino Médio conseguem compreender conceitos básicos de \textit{machine learning} e aplicá-los em tarefas com algoritmos. \citeonline{sanusi2023systematic}, em revisão sobre a Educação Básica, também identificam experiências em diferentes etapas, com presença expressiva do Ensino Médio. As pesquisas apontam lacunas que vão da integração curricular e formação docente até a avaliação das intervenções. Assim, a adequação à faixa etária depende menos da idade e mais do recorte escolhido, das representações adotadas e do suporte oferecido na aula.

Esses estudos afastam a ideia de que o Ensino Médio seja inadequado para discutir modelos simples baseados em dados, embora alertem que o simples uso de recursos digitais está longe de garantir compreensão. As intervenções analisadas variam bastante em duração e objetivos conforme os recursos adotados. Uma experiência conduzida em uma escola específica dificilmente vira receita automática para outra. A literatura sustenta uma proposta viável, cabendo à aplicação revelar seus resultados.

A sequência não deve começar pela expressão matemática de um algoritmo, pois isso pode afastar estudantes que ainda consolidam a leitura de gráficos e noções de proporcionalidade. A manipulação de controles e a repetição de nomes técnicos também podem ocultar o raciocínio que se pretende desenvolver. A introdução ao tema precisa partir de problemas reconhecíveis para que a turma examine os dados e suas representações antes de discutir a regra que gera a saída. Essa progressão cria uma porta de entrada compatível com o nível de escolaridade previsto, sem diminuir o conteúdo de Computação.

O vocabulário técnico não é a única dificuldade esperada. O estudante pode precisar localizar pontos nos dois eixos para comparar distâncias, além de interpretar como as cores e curvas do gráfico expressam a regra do modelo. A leitura da visualização precisa ser ensinada como parte da atividade, pois observar uma mudança não significa explicar a relação que a produziu.

Há também um obstáculo matemático. Noções como média e distância, ou mesmo função e probabilidade, nem sempre estão estabilizadas na turma. A sequência passa longe de resolver, em oito encontros, lacunas acumuladas na escolarização matemática. Ela retomará esses conceitos apenas para responder a perguntas concretas sobre o modelo, mantendo a Matemática presente em um nível de formalização acessível.

Essa retomada não transforma o percurso em revisão extensa de Matemática antes do conteúdo de Computação. A necessidade aparece dentro de cada problema, como na comparação entre previsão e valor observado na regressão linear ou na distância entre exemplos no \textit{k}-NN. A mediação deverá identificar quando uma retomada curta basta e quando a dificuldade impede que a turma prossiga com a interpretação planejada. A seção 3.3 detalha como a Matemática entra nas representações.

Uma interface visual pode reduzir algumas barreiras e, ao mesmo tempo, ocultar outras. Um controle que altera o número de vizinhos torna uma mudança visível, mas não explica, por si só, por que certos pontos passam a decidir uma categoria. O recurso visual precisa ser acompanhado por mediação orientada, como se discute na seção 3.3.

A plataforma trabalha regressão linear, regressão logística, \textit{k}-vizinhos mais próximos, árvore de decisão e \textit{k}-means. A seleção abrange desde a previsão de valores numéricos e a classificação até o agrupamento de dados não rotulados. Em cada módulo, o estudante parte do problema proposto para levantar uma hipótese antes do experimento, explicando depois como a mudança nos dados ou no parâmetro alterou o gráfico e o resultado.

Os modelos foram escolhidos porque mostram mecanismos distintos em representações examináveis por estudantes iniciantes. Na regressão linear, o foco recai sobre a linha de tendência e a relação entre previsão, resíduo e erro. Já a regressão logística separa probabilidade e classificação, pois variar o limiar muda a decisão sem alterar o cálculo de base. Enquanto o \textit{k}-NN classifica pelo confronto com vizinhos semelhantes e a árvore de decisão explicita condições sucessivas, o \textit{k}-means forma grupos mesmo na ausência de categorias prévias para cada ponto.

Esse recorte reúne duas formas distintas de aprendizagem a partir de dados. Nos quatro primeiros módulos, que cobrem das regressões ao \textit{k}-NN e à árvore de decisão, os exemplos já trazem respostas conhecidas para orientar a previsão ou a escolha de categoria. No \textit{k}-means, não há categoria fornecida previamente, e o objetivo é formar grupos a partir de uma medida de proximidade. \citeonline{touretzky2023machine} situam essa diferenciação entre aprendizagem supervisionada e não supervisionada como uma das ideias centrais para a educação em IA. Neste projeto, ela será abordada com exemplos sintéticos e visualizações.

Os cinco módulos não devem ser assimilados com a mesma velocidade ou pelo mesmo caminho. A regressão linear pode se aproximar de conteúdos matemáticos já conhecidos, enquanto o agrupamento sem rótulos pode exigir mais tempo para que os estudantes diferenciem descobrir grupos e separar categorias previamente definidas. A retomada no encontro seguinte e a comparação coletiva no sétimo encontro evitam tratar cada modelo como episódio isolado.

O nível introdutório não exige ocultar os limites dos modelos. Uma reta pode resumir a tendência dos dados observados e ser afetada por pontos atípicos. Previsões fora do intervalo observado exigem cautela, pois a tendência pode não se manter. Uma porcentagem representada na regressão logística não é certeza.

No \textit{k}-NN, por exemplo, a decisão varia conforme os exemplos guardados e a escala adotada, além do número de vizinhos escolhido. Em árvores, o excesso de regras pode falhar diante de casos novos, enquanto os grupos do \textit{k}-means mudam segundo a quantidade e a posição inicial dos centroides. Tratar esses limites na leitura dos gráficos evita tanto a promessa de precisão automática quanto um debate moralizante e abstrato. Ao modificar dados e parâmetros nos controles de cada modelo, o estudante observa a consequência visual e discute qual regra a produziu.

\section{Representações, visualização e experimentação guiada}

Ferramentas visuais permitem introduzir \textit{machine learning} a estudantes iniciantes. \citeonline{gressevonwangenheim2021visual} identificaram ferramentas que favorecem a construção interativa de modelos, embora várias se concentrem na classificação de imagens e em atividades curtas. O produto proposto articula dados sintéticos e tabelas da própria turma em tarefas de regressão e classificação, além de agrupamento. Essa diferença de escopo não resolve, por si só, as limitações apontadas pela literatura.

A interação ganha sentido pedagógico quando o estudante investiga uma relação. A visualização pode tornar visível algo que permaneceria abstrato em uma explicação verbal, mas também pode dar a impressão de que observar uma mudança basta para compreendê-la. Os elementos visuais precisam ter significado explícito, e a regra do modelo deve estar acessível.

Os dados sintéticos permitem controlar o que será comparado e produzir situações em que a modificação de um elemento gera uma consequência perceptível. Não servem como previsão confiável sobre pessoas reais. Os dados organizados pela própria turma no segundo encontro mostram de onde vêm os números que serão interpretados. O cenário didático leva a turma a identificar a mudança ocorrida e a explicar qual elemento do modelo produziu aquele efeito.

Articular tabelas e gráficos a fórmulas, textos curtos e controles interativos mobiliza múltiplas representações. \citeonline{ainsworth2006deft} argumenta que representações diferentes podem cumprir funções complementares durante a aprendizagem, mas exigem tarefas de interpretação e integração. A plataforma não deve apenas reunir recursos na mesma tela. Cada representação precisa responder à pergunta examinada pela turma.

Qualquer ajuste nos controles da tela, seja na inclinação da reta e no limiar de decisão, seja nos parâmetros de vizinhos, condições ou centroides, pede uma questão que leve o estudante a antecipar e ler a mudança. Na regressão linear, os pares da tabela ganham forma no plano cartesiano para que a reta indique tendência ou previsão. Na regressão logística, a curva de probabilidade é lida em conjunto com o limiar que define a classe. Já no \textit{k}-NN, a posição do ponto consultado permite medir a distância até os exemplos para acompanhar a votação dos vizinhos. Enquanto nas árvores o percurso visual reflete as condições testadas, no \textit{k}-means os centroides funcionam como referências móveis que reorganizam os grupos, longe de serem pontos decorativos. Em todos os módulos, a visualização vincula os dados e a regra do modelo à saída gerada.

A atividade adapta a estratégia de prever, observar e explicar proposta por \citeonline{white1992probing}. Antes de usar o controle, o estudante registra uma hipótese oral ou escrita e, na sequência, acompanha como o gráfico altera a medida calculada e a decisão para explicar se o resultado confirmou sua expectativa. O ciclo evita a tentativa e erro em que parâmetros são movimentados sem relação clara entre entrada e saída.

A previsão torna visível uma ideia que poderia permanecer implícita. Se um estudante afirma que aumentar o número de vizinhos não alterará uma classificação e encontra decisão diferente, o professor pode perguntar quais exemplos passaram a ser considerados e por que a votação se alterou. Quando o estudante prevê que mover um centroide não terá efeito e observa a reorganização dos pontos, a explicação pode retomar a noção de proximidade. A hipótese inicial não precisa estar correta para ter valor pedagógico.

A aplicação está prevista para o laboratório de informática da escola, com acesso à internet, sujeito à confirmação das condições do campo. Planeja-se que cada estudante use um computador do laboratório ou, conforme a disponibilidade confirmada, um Chromebook, e a realização das atividades nos encontros busca reduzir as dificuldades decorrentes da restrição de acesso fora da escola. Para estudantes com insegurança no manuseio do computador, os controles deverão ser legíveis e indicar o que pode ser modificado em cada momento. Para aqueles que operam a interface com facilidade, o desafio deixa de ser o manuseio rápido do controle e passa a ser a explicação. A mediação docente, com suporte individual quando necessário, ajuda a evitar que a técnica de uso se torne obstáculo ou substitua a questão conceitual.

O controle interativo não será tratado como decoração ou recompensa após uma explicação pronta. Ele permite comparar uma consequência ao manter alguns fatores estáveis e alterar um dado ou parâmetro. A aula guiada orienta, inicialmente, uma alteração por vez, pois mudanças simultâneas dificultam atribuir uma consequência a uma causa percebida.

Depois da aula guiada, a exploração livre poderá ampliar as comparações. O professor deverá retomar a finalidade do experimento e solicitar explicações que não se limitem a descrever elementos apresentados.

As representações matemáticas serão apresentadas em camadas. O percurso parte de uma situação concreta e do uso de tabelas ou gráficos na leitura dos dados, passando por cálculos simples até chegar à fórmula que resume a regra discutida. Assim, tanto a equação da reta e o limiar quanto o cálculo de distâncias, condições ou médias dos centroides entram para explicar mudanças visíveis na tela, evitando a cobrança de memorização desarticulada.

Apresentar a Matemática em camadas não retira precisão dos conceitos. A organização em camadas define o que o estudante precisa para uma compreensão inicial e o momento em que cada noção passa a fazer sentido. Antes de calcular a distância entre dois pontos, o estudante precisa reconhecer por que a proximidade importa no \textit{k}-NN, e antes de discutir o erro de uma reta, precisa comparar uma previsão com um valor observado.

Introduzida após essas experiências, a fórmula ganha sentido prático por nomear uma relação que o estudante já conseguiu examinar visualmente durante o experimento.

\section{Avaliação da compreensão durante a intervenção}

Avaliar o ensino de \textit{machine learning} na Educação Básica permanece uma questão em aberto. \citeonline{rauber2023assessing} levantaram avaliações que vão de questionários a tarefas baseadas em desempenho e destacaram a necessidade de abordagens mais abrangentes.

\citeonline{gressevonwangenheim2022proposal} apresentam uma rubrica baseada em desempenho para conceitos e práticas de \textit{machine learning}, mostrando a importância de converter expectativas de aprendizagem em critérios observáveis. Embora voltado à classificação de imagens, esse modelo pode ser adaptado ao recorte desta pesquisa. Reconhecer o nome de um algoritmo está longe de provar que o estudante compreende a tarefa resolvida ou a ligação entre dados e saída. Como um acerto numérico pode decorrer de simples chute na leitura do gráfico ou da repetição de falas docentes, os instrumentos pedirão justificativas breves que revelem o raciocínio adotado.

A compreensão vai além do acerto em questões de múltipla escolha. O diagnóstico inicial e o núcleo comparável do instrumento final solicitarão operações cognitivas equivalentes em extensão e dificuldade. Enquanto registros individuais no segundo, quarto e sexto encontros acompanharão as previsões e explicações da turma, as notas do professor-pesquisador situarão as dúvidas surgidas e as intervenções nas condições de uso.

Longe de selecionar estudantes ou medir déficits, o diagnóstico inicial levantará concepções prévias. Isso ocorre quando alguém atribui decisões à inteligência artificial ignorando dados e regras, ou quando confunde previsão com classificação e supõe categorias prontas em agrupamentos. No instrumento final, o primeiro bloco retomará essas dimensões em questões equivalentes. Já o segundo bloco trará uma situação inédita para examinar a transferência conceitual, evitando equiparar essa tarefa mais complexa ao levantamento feito no início.

\citeonline{black1998assessment} discutem a avaliação formativa como parte do ensino, pois as informações produzidas pelos estudantes podem orientar decisões docentes. As perguntas na plataforma servirão para dar devolutiva imediata e organizar a conversa em sala, sem coleta digital de respostas. A análise confrontará o diagnóstico e o instrumento final com registros individuais selecionados e anotações docentes.

Durante a exploração, o professor poderá usar as previsões para identificar dificuldades que peçam retomada, dispensando o registro de cada fala como dado de pesquisa. Caso vários estudantes tomem uma probabilidade de 70\% como garantia de ocorrência, a aula poderá pausar a manipulação para distinguir estimativa de certeza. Da mesma forma, eventuais confusões entre \textit{k}-NN e \textit{k}-means serão tratadas retomando a presença ou ausência de rótulos nos exemplos. Tais intervenções pertencem ao ensino, cabendo ao registro posterior contextualizar as produções.

A análise se apoia em quatro dimensões que vão do reconhecimento da origem dos dados e da tarefa até a antecipação de ajustes, passando pela leitura de representações e dos limites na decisão automatizada. Uma rubrica em três níveis separará as respostas sem relação justificável daquelas com relações parciais ou explicações consistentes com o caráter introdutório da sequência, organizando o corpus sem gerar um índice isolado de competência.

Uma situação sobre consumo fictício de energia, por exemplo, poderá exigir que o estudante diferencie temperatura e consumo, evitando tratar todos os valores apenas como dados genéricos. Outras questões pedirão que ele antecipe o que acontece ao mover pontos e centroides na tela, ou ao mudar parâmetros como o limiar, as condições e o número de vizinhos. Assim, a leitura de tabelas e gráficos nos modelos será avaliada pela justificativa dada ao acerto.

Há diferença entre escrever que ``mudou porque o programa mudou'' e explicar que a alteração de limiar passou a classificar como alerta valores antes abaixo da regra estabelecida. A rubrica registrará essa diferença por dimensão, acompanhada de justificativa curta, sem produzir rótulos fixos para caracterizar estudantes.

A apresentação dos resultados combinará sínteses por dimensão e tabelas de transição entre momentos comparáveis, apoiadas por frequências descritivas com seus denominadores e breves excertos anonimizados. A análise acompanhará o que muda ou permanece em cada participante, evitando ranquear os estudantes. Já o desafio de transferência e o questionário de percepção terão seções próprias, situando a discussão nas condições que delimitam a intervenção.

\section{Síntese do referencial e posicionamento da proposta}

A literatura e os documentos curriculares definem o posicionamento da proposta. Nela, o \textit{machine learning} entra como conteúdo introdutório de Computação, superando o simples treinamento no uso de aplicativos. A sequência recorta ideias fundamentais e limites observáveis nas tarefas, sem antecipar uma formação universitária em algoritmos. Tanto os dados da própria turma quanto os cenários sintéticos visam mostrar a origem dos números e o efeito de alterar parâmetros. Dessa forma, a visualização sustenta comparações nas quais o estudante parte de uma hipótese para observar o gráfico e explicar como a regra do modelo transforma os dados em resultado.

A análise buscará evidências em produções individuais contextualizadas, indo além da contagem de acertos sem prometer efeitos causais generalizáveis.

Embora a literatura já reúna ferramentas e experiências no ensino de IA e \textit{machine learning} na Educação Básica, os estudos cobram maior clareza sobre conteúdos e estratégias nos contextos de aplicação, ao lado de avaliações mais rigorosas (\MakeUppercase{MARTINS; GRESSE VON WANGENHEIM}, \citeyear{martins2023findings}; \MakeUppercase{RAUBER; GRESSE VON WANGENHEIM}, \citeyear{rauber2023assessing}; \MakeUppercase{SANUSI} et al., \citeyear{sanusi2023systematic}). Assim, este projeto contribui ao analisar, em uma turma de Ensino Médio do próprio pesquisador, uma sequência didática que articula a visualização de cinco modelos a instrumentos com critérios explícitos.

A qualidade da pesquisa dependerá do alinhamento metodológico entre o referencial e o produto desenvolvido. Para isso, os recursos do site refletirão objetivos de aprendizagem presentes nas atividades e nos instrumentos de coleta, garantindo que as conclusões respeitem estritamente o alcance das evidências produzidas.

\chapter{Produto Educacional e Sequência Didática}

\section{Concepção do produto educacional}

O produto educacional une uma plataforma interativa em português brasileiro a uma sequência didática para o Ensino Médio. A plataforma roda direto no navegador com cenários sintéticos, executando todos os cálculos no próprio computador. Seu acesso dispensa criação de conta ou cadastro de usuário, operando sem envio de respostas ou coleta de dados pessoais. Isso facilita o uso mesmo sob conectividade limitada, bastando abrir o recurso nos computadores ou no projetor. A ferramenta apoia a mediação docente em Computação, preservando o espaço da explicação e do diálogo conduzidos pelo professor.

A proposta tem cinco módulos introdutórios de \textit{machine learning}: regressão linear, regressão logística, \textit{k}-vizinhos mais próximos, árvore de decisão e \textit{k}-means. Longe de formarem um catálogo de comandos ou conteúdos soltos, eles apresentam tarefas distintas de previsão numérica e classificação, além da formação de grupos para que a turma compare o que muda entre eles. Todos ocupam lugar central na sequência, mobilizando camadas matemáticas próprias. Na regressão linear, o estudo da função e da linha de tendência fundamenta a previsão e o cálculo de resíduos ou erros, enquanto a logística articula probabilidade ao limiar de decisão. Já o \textit{k}-NN destaca a distância entre pontos, o \textit{k}-means recalcula centroides pela média e a árvore encadeia condições.

A diferença de formalização não diminui nenhum módulo. Ela decorre do repertório de estudantes iniciantes e da necessidade de tornar visível, no gráfico ou na regra, uma relação que a turma consiga discutir durante a aula.

Mais do que estarem na mesma plataforma, os modelos se conectam por perguntas recorrentes nos módulos. A turma identifica a origem do dado disponível e a saída esperada antes de testar qual ajuste pode alterá-la. A resposta varia em cada caso. Enquanto a reta gera uma previsão numérica, a regressão logística calcula uma probabilidade confrontada com um limiar. Por sua vez, o \textit{k}-NN consulta exemplos próximos e a árvore percorre regras condicionais, ao passo que o \textit{k}-means reorganiza pontos mesmo sem categorias prévias.

Com as mesmas perguntas em cada módulo, os estudantes comparam os modelos sem reduzi-los a uma lista de nomes. O professor percebe se a dificuldade está em um conceito específico ou na leitura de dados e gráficos, que atravessa todos os módulos.

\section{Organização pedagógica da plataforma}

Cada módulo parte de uma situação-problema e da origem dos dados para questionar o que o modelo tenta aprender. Antes de usar os controles, o estudante prevê uma consequência. Em seguida, modifica dados ou parâmetros definidos para o experimento, observando o efeito gerado no gráfico e na saída do modelo para explicar o que ocorreu. Esse movimento da previsão ao experimento, seguido pela observação e pela explicação, repete-se nos módulos para explicitar o raciocínio por trás de cada manipulação.

Mover a inclinação da reta leva o estudante a comparar a previsão com o valor observado, enquanto a troca de um limiar questiona por que a classe muda sem alterar a porcentagem calculada. A mesma lógica orienta o ajuste de vizinhos ou de divisões da árvore, assim como a posição dos centroides.

A sequência abre com recursos de inteligência artificial que os estudantes já usam e chega logo a um problema concreto com dados. A Matemática entra à medida que a leitura do gráfico e a decisão do modelo na situação exigem. Toda representação visual é acompanhada por tabelas e textos que explicitam a regra, evitando que a compreensão dependa apenas da imagem. As visualizações em 3D do protótipo entram como recurso complementar dispensável ao conceito central, mantendo na interface o uso por teclado e a opção de reduzir movimentos.

\section{Cenários e dados dos experimentos}

Os dados das atividades serão quase todos sintéticos. No segundo encontro, a turma poderá montar uma pequena tabela anônima com informações simples e voluntárias, ou mesmo medir objetos da sala. Essa tabela servirá apenas ao ensino, ficando fora do corpus. Os demais cenários descartam o uso de notas escolares, registros de comportamento ou qualquer dado pessoal dos estudantes.

Os cenários constituem uma proposta inicial sujeita a ajustes após a definição da escola e da turma. Para a regressão linear, a ideia é relacionar temperatura e consumo fictício de energia em uma sala, enquanto a regressão logística usa a leitura de um sensor simulado para discutir probabilidade e emissão de alerta. Já o \textit{k}-NN classificará situações descritas por duas medidas, ao passo que a árvore de decisão simulará o acionamento de uma irrigação conforme a umidade e a temperatura. No \textit{k}-means, pontos de coleta em um mapa fictício de bairro permitirão formar grupos não rotulados a partir da proximidade com os centroides. Assim, além do apelo visual, cada contexto deve tornar visível o efeito de alterar dados ou parâmetros em um nível compatível com a Matemática da série. Antes da aplicação, os exemplos serão revistos para adequar a linguagem e o tempo previsto ao perfil da turma.

\section{Sequência didática provisória}

A sequência foi planejada, por ora, para oito encontros de cerca de cinquenta minutos. No campo inicialmente previsto, o laboratório de informática da escola já conta com 20 computadores conectados à internet; caso a aplicação mude de unidade, a infraestrutura local precisará ser conferida. A confirmação da turma e do número de aulas dentro do componente curricular ainda depende da viabilidade do campo. O cenário preferencial é uma turma de Ensino Médio do próprio pesquisador, possivelmente na disciplina Projeto Permanente por Afinidade, podendo migrar para Matemática ou outro espaço curricular conforme a lotação e o calendário exigirem.

Todo encontro mantém a mesma estrutura, mesmo quando o modelo muda. A aula abre com uma situação ou pergunta que delimita o problema. Após registrar uma breve previsão individual, cada estudante explora os controles individualmente antes do debate coletivo e da síntese docente no fechamento. Essa rotina dá continuidade às aulas e integra a plataforma ao trabalho de construir e revisar hipóteses, evitando que o uso do site funcione como simples intervalo entre explicações.

O primeiro encontro serve para nivelar conceitos. Antes de qualquer modelo, a turma discute aplicações cotidianas de inteligência artificial. Depois distingue os dados e a regra do modelo, até apontar o que seria entregue ao sistema e o que se espera dele. O diagnóstico inicial é individual e acontece no começo desse encontro, antes da discussão, ou na aula anterior, conforme o tempo. Registra ideias de partida, sem antecipar os nomes técnicos que aparecem depois.

Do segundo ao sexto encontro entram, um por aula, os cinco módulos da plataforma, indo das regressões aos modelos de classificação e agrupamento. No segundo encontro, a organização da tabela da turma antecede a regressão linear. Cada aula retoma elementos anteriores para evitar que os modelos fiquem isolados em uma única experiência.

A regressão linear e a regressão logística abrem a parte dos modelos supervisionados. Na linear, a turma identifica a linha de tendência no plano cartesiano para fazer previsões e compará-las aos valores observados, introduzindo a noção de erro. A logística retoma o uso de exemplos conhecidos e direciona o debate para o cálculo de probabilidade e a escolha de classe pelo limiar. Essa sequência imediata permite contrastar a previsão de um valor numérico com a definição de uma categoria, mostrando que ambas dependem dos dados e das escolhas do modelo.

O \textit{k}-NN e a árvore de decisão são duas maneiras de chegar a uma classificação. No \textit{k}-NN, os estudantes observam o papel da distância e do número de vizinhos na votação; na árvore, acompanham regras condicionais e os caminhos até uma decisão. Comparar os dois afasta a ideia de que exista um algoritmo melhor que os outros. Ambos usam dados e regras, cada um à sua maneira, e entregam saídas que precisam ser lidas à luz do problema e dos exemplos disponíveis.

O \textit{k}-means traz uma mudança: os pontos não vêm com categorias dadas de antemão. A atividade volta à pergunta sobre o que se quer produzir com os dados e mostra que formar grupos por proximidade não é escolher uma classe conhecida. A atualização dos centroides e a observação da reorganização dos pontos poderão exigir mais tempo, pois a diferença entre agrupamento e classificação pode levar o professor a retomar a linguagem dos encontros anteriores. Uma síntese breve registra o modo como os cinco modelos articulam dados e regras para gerar uma saída, destacando o que muda conforme a finalidade de cada tarefa.

O sétimo encontro compara os cinco modelos. Diante de novas situações-problema, a turma justifica com base nos dados disponíveis se o caso exige previsão numérica, escolha de categoria ou formação de grupos. No oitavo encontro, cada estudante responde primeiro ao bloco de questões equivalentes às do diagnóstico e, na sequência, resolve um desafio de transferência inédito. Nessa etapa final, é preciso apontar os dados necessários para associar a tarefa a um modelo compatível, explicando o efeito esperado caso um parâmetro ou dado mude na regra.

O encontro comparativo recupera diferenças que possam ter ficado dispersas e prepara a transferência para o desafio final. Ao justificar que tipo de saída o problema demanda a partir dos dados disponíveis, o estudante considera a finalidade da tarefa antes de escolher o modelo. Enquanto o núcleo comparável permite examinar mudanças frente ao diagnóstico, o desafio de transferência mostra como os conceitos são mobilizados em novo contexto, entrando como evidência adicional sem equiparação direta com questões mais simples.

O trabalho na plataforma ocorrerá individualmente, com os computadores do laboratório e, conforme a disponibilidade confirmada no campo, Chromebooks para completar o acesso da turma. No segundo, no quarto e no sexto encontros, cada estudante também preencherá um registro curto antes e depois de explorar o site, cobrindo assim um exemplo de regressão, um de classificação e outro de agrupamento. Já o instrumento final será inteiramente individual. Caso o laboratório fique indisponível, o professor projetará os experimentos para a turma, preservando as pausas de previsão e explicação individual.

O roteiro preliminar dos oito encontros articula objetivos de aprendizagem e tempo previsto às ações de mediação com questões formativas, indicando ainda adaptações de infraestrutura e limites didáticos por módulo. Sem armazenamento de respostas digitadas no site, o corpus reunirá o diagnóstico inicial e o instrumento final completo, somados aos três registros individuais intermediários e às anotações docentes.

A versão do site usada no campo será conferida, antes da intervenção, com os recursos previstos na sequência. Se o cenário da plataforma tiver dados diferentes dos discutidos no roteiro da aula, a diferença será corrigida ou explicada aos estudantes; o mesmo vale para os textos de apoio e as representações matemáticas.

A pesquisa não pode analisar estudantes como se tivessem passado por uma sequência diferente do recurso usado. Essa conferência prática garante que a plataforma e os instrumentos avaliem exatamente o que foi trabalhado nas aulas.

\section{Materiais e limites do produto}

Além da plataforma, a aplicação contará com o diagnóstico inicial e três registros curtos de aula, somados às notas docentes e ao instrumento final com seu questionário de percepção. Nos registros intermediários, um campo recolhe a previsão prévia ao experimento e outro a explicação posterior. Já o questionário final consulta a turma sobre a clareza e o nível de dificuldade no manuseio do site, sem valer como evidência de aprendizagem.

O produto não comprova aprendizagem só por existir. Serve para que professor e estudantes discutam dados e parâmetros do modelo em sala.

A adequação didática do produto só poderá ser avaliada confrontando o planejamento com as condições reais do campo e as produções da turma, longe de constituir solução pronta para qualquer escola. Antes das aulas, um teste prático conferirá a abertura local do site nos computadores e a legibilidade da tela, checando ainda a resposta dos controles e a visibilidade no projetor. Caso algum equipamento falhe, o experimento seguirá projetado com pausas para registro individual.

A mudança de organização vai para o registro docente, porque a forma de acesso ao recurso pode alterar as interações observadas. Manter uma atividade pedagogicamente equivalente vale mais do que insistir em uma configuração técnica que a escola não consiga sustentar.

A mesma verificação confere se os estudantes compreendem os enunciados sem domínio prévio do jargão da área. Termos como característica e limiar, ou mesmo vizinho e centroide, entram sempre acompanhados de exemplos visíveis e explicações situadas. Caso alguma pergunta gere dúvida recorrente, a revisão do enunciado ou da mediação será tratada como ajuste necessário ao produto, afastando a ideia de falha do estudante para remover barreiras linguísticas alheias ao conteúdo.

\section{Alinhamento curricular preliminar}

Uma matriz explicitará o vínculo da sequência com a Computação na Educação Básica, associando cada habilidade curricular ao seu objetivo de aprendizagem, bem como à atividade proposta e às evidências observáveis. A versão preliminar apoia-se nas referências do Ensino Médio voltadas aos fundamentos da inteligência artificial com suas potencialidades e limites, articulando a análise em ciência de dados ao uso de modelos computacionais simples de simulação e previsão. Esse recorte dialoga com as habilidades EM13CO10 a EM13CO12 previstas no anexo da Resolução CNE/CEB nº 1/2022 \cite{brasil2022computacao}. Embora a norma deixe de prescrever esses cinco algoritmos nominalmente, eles oferecem situações didáticas concretas para examinar os limites de cada regra na transformação de dados em resultados. A matriz indicará ainda conexões com a Matemática, mantendo a Computação como objeto principal. O recorte definitivo das habilidades dependerá da série e da disciplina definidas, respeitando o currículo da rede estadual e as condições do campo.

\chapter{Percurso Metodológico}

\section{Natureza, abordagem e estratégia da pesquisa}

A pesquisa será de natureza aplicada, com abordagem predominantemente qualitativa e apoio de dados quantitativos descritivos. Será desenvolvida como intervenção pedagógica de grupo único, com produções individuais em diferentes momentos da sequência. Esse desenho permite examinar mudanças, permanências e dificuldades nas explicações dos estudantes, sem pretender isolar um efeito próprio da plataforma frente à mediação docente, às atividades e às interações em sala.

A estratégia principal será uma intervenção pedagógica. \citeonline{damiani2013discutindo} caracterizam esse tipo de investigação pelo planejamento e pela implementação de mudanças ou inovações pedagógicas, acompanhados da análise de seus efeitos. Essa formulação é compatível com o projeto porque os recursos da intervenção serão planejados para uma situação de ensino, aplicados no campo e examinados a partir das evidências produzidas.

A intervenção ganha sentido como uma unidade integrada que articula a sequência didática e a plataforma aos instrumentos de produção de dados e à mediação docente. Separar artificialmente esses elementos seria inadequado, pois, nesta pesquisa, a plataforma adquire o sentido previsto em conjunto com a pergunta proposta, a previsão solicitada e a discussão posterior.

Reconhecer essa unidade não autoriza atribuir qualquer mudança observada a um único componente. As produções dos estudantes entram como evidências construídas em um percurso situado, no qual a infraestrutura, o tempo de aula, o acesso individual à plataforma e as retomadas conceituais moldam a forma como cada explicação aparece. A posição do professor-pesquisador exigirá registro claro das condições de aplicação e prudência na interpretação dos dados.

A definição metodológica final deverá ser discutida com o orientador antes da submissão do projeto e do início da produção de dados. Essa decisão dependerá também das condições confirmadas do campo e das exigências éticas e institucionais aplicáveis.

\section{Contexto e participantes}

O campo previsto é uma turma de Ensino Médio de escola pública na qual o pesquisador atue como professor. A possibilidade inicial é uma turma vinculada à disciplina Projeto Permanente por Afinidade, em escola de tempo integral. Esse cenário não está garantido para o próximo ano letivo, pois a lotação docente e a composição das turmas poderão mudar.

Se necessário, a sequência poderá ser aplicada em outra turma do Ensino Médio sob docência do pesquisador, inclusive em um componente de Matemática, desde que sejam preservadas as condições pedagógicas e éticas definidas para a pesquisa.

Os participantes serão estudantes iniciantes em \textit{machine learning}, sem exigência de programação prévia. A sequência pressupõe apenas que saibam ler tabelas e gráficos elementares, localizar pontos em eixos cartesianos e comparar valores, além de mobilizar noções iniciais de função e probabilidade, ou de distância e média, conforme forem retomadas nos módulos.

A definição da escola e da série, junto ao número de participantes no componente curricular e ao calendário, ocorrerá após confirmar o campo.

\section{Etapas da intervenção}

O primeiro momento corresponderá ao estudo das referências para aperfeiçoar os exemplos sintéticos da plataforma e os instrumentos previstos. Em seguida, a sequência didática será preparada de acordo com a turma e as condições de infraestrutura disponíveis.

Prevê-se que os estudantes trabalhem individualmente, com um computador do laboratório ou, conforme a disponibilidade confirmada, um Chromebook. No campo inicialmente previsto, o laboratório dispõe de 20 máquinas; em outro campo, será necessário verificar as condições disponíveis. O acesso individual permite que cada estudante formule e teste suas próprias hipóteses diante do mesmo experimento. As previsões e explicações de interesse serão anotadas individualmente antes e depois da exploração em três encontros, enquanto os registros do professor indicarão de que modo o acesso individual, a infraestrutura ou as retomadas conceituais ajudam a dimensionar essas produções. Os oito encontros asseguram uma duração viável, exigindo atenção constante ao ritmo da turma.

Ao final de cada encontro, o professor-pesquisador organizará uma síntese breve que recupere os elementos discutidos pela turma.

As sínteses de encerramento não terão a função de produzir uma resposta correta única para ser copiada pelos estudantes. Elas servirão para retomar o que foi observado e clarear relações confusas antes da passagem ao encontro seguinte. Caso a turma confunda uma probabilidade com garantia de ocorrência, por exemplo, a retomada poderá recuperar o papel do limiar e distinguir estimativa de certeza antes de avançar para outro modelo.

O registro docente indicará que a retomada ocorreu e quais dificuldades a motivaram, sem substituir as respostas individuais como fonte de evidência da compreensão. Essa separação ajuda a preservar o vínculo entre a mediação necessária ao ensino e a análise cuidadosa do que cada participante conseguiu elaborar.

A plataforma não registrará respostas ou dados pessoais. As respostas usadas na pesquisa serão produzidas nos instrumentos descritos a seguir.

\section{Instrumentos de produção de dados}

Antes da sequência, os participantes realizarão um diagnóstico individual breve. O instrumento trará um núcleo de quatro ou cinco questões para examinar desde os dados e a finalidade da tarefa até a interpretação de representações, o efeito de alterações e os limites de decisões automatizadas. Ele servirá para registrar o ponto de partida das explicações e não comporá a avaliação escolar.

Nos encontros 2, 4 e 6, cada participante preencherá um registro curto com duas solicitações: indicar o que prevê antes de alterar o controle e explicar o resultado depois da exploração. A escolha desses encontros permite acompanhar uma tarefa de previsão numérica, uma tarefa de classificação e uma tarefa de agrupamento, sem exigir produção escrita em todos os módulos.

O instrumento final individual terá dois blocos. O primeiro retomará o núcleo do diagnóstico por meio de questões equivalentes em operação cognitiva, extensão e dificuldade. Os contextos poderão mudar, mas cada par de questões deverá solicitar a mesma relação conceitual. O segundo bloco será um desafio de transferência. Nele, o estudante analisará uma situação inédita para reconhecer a tarefa e os dados necessários, indicando um modelo compatível e explicando o efeito esperado caso um dado, parâmetro ou regra sofra alterações.

O desafio de transferência não será tratado como equivalente ao diagnóstico. Sua função será examinar a mobilização dos conceitos em outro contexto. Um questionário breve de percepção registrará a utilidade atribuída às atividades e à plataforma, bem como a clareza e o nível de dificuldade encontrados. Essas respostas serão analisadas como apreciação do produto, e não como demonstração de aprendizagem.

Após cada aula, o professor-pesquisador estruturará anotações em tópicos cobrindo as intervenções realizadas e as dificuldades recorrentes, além do tempo empregado e das condições de uso observadas. As perguntas inseridas no site terão função formativa imediata. Como suas respostas não permanecem salvas, não serão usadas diretamente como dados da pesquisa.

Antes da aplicação, uma matriz associará cada questão comparável ao objetivo de aprendizagem correspondente e à atividade desenvolvida, ligando a dimensão da rubrica ao item equivalente no outro instrumento. Os instrumentos serão submetidos à leitura crítica do orientador e de um professor com experiência no nível de ensino. Quando for viável, a clareza dos enunciados será verificada com pessoa de perfil próximo ao público que não participe da pesquisa.

Os instrumentos também precisarão ser compatíveis com o tempo escolar disponível e com o fato de que os participantes serão estudantes do Ensino Médio em uma atividade inserida no cotidiano da turma. Não se pretende aplicar teste extenso que interrompa a sequência ou transforme a intervenção em rotina de avaliação formal.

Perguntas abertas curtas, acompanhadas de gráfico, tabela ou situação-problema, podem oferecer indícios mais úteis que grande quantidade de itens objetivos, desde que o enunciado deixe claro o que se pede. Antes da aplicação, o professor-pesquisador verificará se cada questão está ligada a um objetivo de aprendizagem e a uma atividade efetivamente realizada.

Caso algum conceito não tenha sido abordado nas condições previstas, ele não deverá ser tratado como ausência de compreensão no instrumento final. Essa precaução mantém os instrumentos vinculados ao percurso de ensino e protege a análise contra interpretações que ultrapassem as evidências disponíveis.

\section{Procedimentos de análise}

As produções serão organizadas por códigos, sem identificação nominal dos participantes. A rubrica terá três níveis, variando de 0 a 2, estruturada em quatro dimensões que vão dos dados e da tarefa até a representação conceitual e os limites da decisão. O nível 0 reunirá respostas ausentes, incompatíveis ou sem relação justificável com a situação. O nível 1 indicará relação parcial entre os elementos. O nível 2 indicará uma explicação consistente com o nível introdutório definido para a sequência. Cada classificação será acompanhada por uma justificativa curta.

A primeira etapa consistirá na leitura integral das respostas e na revisão dos exemplos que delimitam os três níveis. Em seguida, estruturar-se-á uma matriz que cruza cada participante e dimensão aos momentos da coleta. Os códigos apoiarão a organização do corpus, sem virar nota para a turma.

Uma resposta poderá apenas indicar que ``a classificação mudou'', enquanto outra relacionará a mudança ao limiar escolhido ou à proximidade entre exemplos. A diferença entre esses modos de explicar será mais informativa que uma oposição simples entre certo e errado.

As questões comparáveis do diagnóstico e do instrumento final serão colocadas em relação. A análise acompanhará se houve manutenção do nível inicial ou mudança para cima e para baixo, sempre lendo a justificativa escrita. Uma tabela de transição por dimensão mostrará quantos participantes seguiram cada percurso, enquanto os registros intermediários situarão o surgimento dessas relações nas aulas. O desafio de transferência será analisado à parte nas dimensões aplicáveis. Frequências absolutas e relativas virão com seus denominadores. Casos de respostas em branco ou de participantes sem par completo entre os dois instrumentos serão informados sem imputação de valores. Essas medidas descreverão o conjunto e não sustentarão inferências estatísticas ou generalizações para outras turmas.

O texto final trará sínteses por dimensão e frequências descritivas acompanhadas de tabelas de transição, além de breves excertos anonimizados que exemplifiquem formas distintas de explicação. Os registros docentes contextualizarão as produções escritas sempre que dúvidas recorrentes e retomadas matemáticas, ou mesmo condições de infraestrutura, interferirem na atividade. O questionário de percepção será apresentado em seção própria e não será combinado com os níveis da rubrica.

A interpretação considerará as condições em que a atividade ocorreu, mas não usará registros do professor para completar uma resposta individual. Quando for viável, um segundo avaliador classificará uma parcela das produções. As divergências serão discutidas com base nas definições e nos exemplos da rubrica, e os ajustes de critério serão registrados antes da análise final.

\section{Aspectos éticos}

A coleta de dados só terá início após a liberação institucional exigida pelos comitês de ética aplicáveis ao projeto. A Resolução CNS nº 510/2016 dispõe sobre pesquisas em Ciências Humanas e Sociais que utilizam dados obtidos diretamente de participantes ou informações identificáveis, enquanto a Resolução CNS nº 466/2012 estabelece diretrizes gerais para pesquisas envolvendo seres humanos \cite{brasil2016cns510,brasil2012cns466}.

Como a intervenção poderá envolver adolescentes, os procedimentos éticos deverão ser definidos conforme a orientação institucional e a tramitação exigida antes da aplicação.

As respostas receberão códigos de anonimização, mantendo eventuais documentos com dados identificáveis em arquivos separados dos instrumentos destinados à análise. Dados brutos não serão inseridos no código da plataforma nem publicados em repositório aberto. A participação será voluntária, garantindo aos estudantes que optarem por não integrar a pesquisa o acesso a atividades pedagógicas equivalentes sem nenhum prejuízo escolar.

\chapter{Contribuições Esperadas e Limites da Proposta}

\section{Contribuições esperadas}

Espera-se produzir uma proposta de ensino introdutório de \textit{machine learning} situada no campo da Computação na Educação Básica. O trabalho articula uma plataforma local a uma sequência de oito encontros com orientações de mediação docente, além de instrumentos individuais analisados por critérios próprios. Esse conjunto permitirá examinar como os estudantes ligam os dados às regras e aos parâmetros na saída do modelo, explicando seus limites.

A proposta oferecerá aos estudantes iniciantes um percurso que dispensa programação prévia ou formalização matemática avançada. A expectativa é criar condições para que a turma levante hipóteses antes de testar alterações e explicar o efeito observado nas tabelas e nos gráficos. Mais do que atestar domínio técnico profundo dos modelos, a pesquisa busca identificar compreensões conceituais situadas ao longo da intervenção.

A matriz curricular construída para o projeto explicitará o vínculo dos objetivos de aprendizagem e das atividades propostas com as habilidades de Computação do Ensino Médio que articulam análise de dados e modelos simples aos fundamentos de inteligência artificial. Esse desenho busca contribuir para o debate sobre a Computação na Educação Básica, sem supor que a BNCC prescreva algoritmos específicos ou uma organização curricular única.

A investigação também poderá sistematizar uma experiência docente adaptável a outras turmas de Ensino Médio, respeitando a infraestrutura disponível e a mediação necessária frente ao repertório e ao currículo de cada escola. Longe de constituírem solução pronta ou universal, a plataforma e os materiais oferecem subsídios a professores interessados em abordar inteligência artificial e dados para além do uso passivo de ferramentas prontas.

O produto também tornará visível o planejamento que sustenta sua utilização, já que a simples oferta de telas com gráficos e controles seria insuficiente. Por isso, tanto os cenários da sequência quanto as questões de previsão e os critérios de análise indicarão a compreensão esperada e as explicações observáveis em cada atividade.

Essa documentação oferecerá aos professores uma referência para justificar adaptações da proposta, evitando padronizar aulas em contextos distintos. Nessas adaptações, caberá preservar o vínculo do conteúdo de Computação com a Matemática mobilizada sob mediação docente, impedindo que o trabalho se reduza à simples demonstração de ferramentas.

\section{Limites e riscos reconhecidos}

O número de encontros e a duração prevista da intervenção limitam a profundidade com que cada modelo poderá ser abordado. Trabalhar cinco modelos introdutórios em oito aulas exige retomadas e comparações constantes, mas não permite supor que uma única exploração resulte em aprendizagem consolidada. O projeto deverá priorizar conceitos e relações fundamentais, deixando de fora formalizações e procedimentos técnicos que ultrapassem o nível introdutório definido.

A participação dos estudantes pode refletir diferenças tanto no repertório matemático quanto na familiaridade e no acesso aos computadores. Embora a sequência preveja estratégias de acesso e mediação, tais cuidados estão longe de eliminar essas desigualdades prévias. Os registros docentes ajudarão a ler como essas condições interferem nas produções individuais, evitando tomar o modelo ou a plataforma como explicação exclusiva para as respostas observadas.

O estudo ocorrerá em uma turma sob docência do próprio pesquisador, ainda a ser definida. Assim, as evidências produzidas refletirão a mediação realizada naquele contexto específico, considerando tanto os dados sintéticos escolhidos quanto as condições materiais do campo. Afasta-se a pretensão de generalizar resultados para toda a rede pública ou de medir um efeito causal isolado da plataforma.

Comparar o núcleo inicial ao final permitirá examinar mudanças e permanências nas explicações, enquanto o desafio de transferência trará evidências sobre o uso dos conceitos em nova situação. Ainda assim, tais achados são insuficientes para atribuir à plataforma a causa isolada das mudanças.

Há limites próprios dos modelos e das representações adotadas. Os conjuntos de dados serão simplificados para evidenciar relações, evitando confundir as saídas da interface com previsões confiáveis sobre pessoas ou situações reais. Além disso, as vistas tridimensionais exibirão apenas métricas calculadas, sem acrescentar novos atributos aos dados.

Os elementos didáticos serão revistos antes da aplicação para afastar a ideia de neutralidade ou infalibilidade em sistemas tidos como autônomos na inteligência artificial.

\chapter{Cronograma Provisório}

O cronograma será ajustado após a confirmação da lotação docente com a turma e o componente curricular, respeitando o calendário escolar e os trâmites éticos institucionais. A Tabela~\ref{tab:cronograma} organiza a sequência de trabalho prevista, subordinada aos prazos oficiais do programa e à definição prática do campo.

\begin{table}[htbp]
\centering
\caption{Etapas e atividades do cronograma provisório da pesquisa}
\label{tab:cronograma}
\begin{tblr}{
  width = \textwidth,
  colspec = {X[1.2,l,m] X[3.8,l,m] X[3.0,l,m]},
  hline{1,Z} = {1.5pt},
  hline{2} = {1pt},
  rows = {abovesep=4pt,belowsep=4pt},
}
\textbf{Etapa} & \textbf{Atividades principais} & \textbf{Condição para conclusão} \\
1. Consolidação do projeto & Revisar o texto de qualificação, alinhando a sequência e a matriz curricular aos instrumentos. & Versão de qualificação coerente e submetida à apreciação do orientador. \\
2. Aperfeiçoamento do produto & Adequar cenários sintéticos e controles do site às aulas, testando a acessibilidade e a projeção. & Protótipo compatível com a aplicação planejada. \\
3. Preparação do campo & Confirmar a escola e a turma no componente curricular, checando a infraestrutura e as autorizações. & Campo definido e condições pedagógicas e éticas confirmadas. \\
4. Revisão dos instrumentos & Ajustar a rubrica aos instrumentos inicial e final, incluindo os registros de aula e de percepção. & Instrumentos revisados e prontos para uso. \\
5. Aplicação da intervenção & Realizar os oito encontros e produzir as respostas individuais e os registros docentes previstos. & Produções organizadas por códigos, sem identificação nominal no corpus de análise. \\
6. Organização e análise & Anonimizar e analisar o corpus com a rubrica, gerando as tabelas de transição e a seleção de excertos. & Mudanças, permanências e dificuldades interpretadas por dimensão. \\
7. Redação da dissertação e do produto & Redigir a análise dos resultados e as considerações finais, fechando a versão definitiva do produto. & Versão submetida às revisões do orientador. \\
8. Revisão e entrega & Conferir a normalização acadêmica e os requisitos éticos ou de acessibilidade para o depósito final. & Trabalho e produto preparados para a etapa institucional correspondente. \\
\end{tblr}
\source{Elaborado pelo autor (2026).}
\end{table}


\phantompart
		
% ELEMENTOS PÓS-TEXTUAIS
\postextual

% Referências bibliográficas
\bibliography{referencias/ref_documento}

\phantompart
\printindex

\end{document}
